% Slug: visual-lexeme-spans
% Source: figures/lexeme-spans.tex:3 (tikzpicture env start line)
% Original label: \label{fig:visual-lexeme-spans}
% Original caption (LaTeX source):
%   Half-open source spans distinguish three tokens in five characters; the final cursor points just past the input.
% Original caption (plain text):
%   Half-open source spans distinguish three tokens in five characters; the final cursor points just past the input.
% Required TikZ setup: \usepackage{tikz} + \usetikzlibrary{arrows.meta, positioning, shapes.geometric, fit} (loaded by support/compiler-pseudocode.sty).
% Shared styles used in this picture: `cfg edge'.
%   Defined in support/compiler-pseudocode.sty: `basic block', `cfg edge', `back edge', `block label'.
% Usage: % Slug: visual-lexeme-spans
% Source: figures/lexeme-spans.tex:3 (tikzpicture env start line)
% Original label: \label{fig:visual-lexeme-spans}
% Original caption (LaTeX source):
%   Half-open source spans distinguish three tokens in five characters; the final cursor points just past the input.
% Original caption (plain text):
%   Half-open source spans distinguish three tokens in five characters; the final cursor points just past the input.
% Required TikZ setup: \usepackage{tikz} + \usetikzlibrary{arrows.meta, positioning, shapes.geometric, fit} (loaded by support/compiler-pseudocode.sty).
% Shared styles used in this picture: `cfg edge'.
%   Defined in support/compiler-pseudocode.sty: `basic block', `cfg edge', `back edge', `block label'.
% Usage: % Slug: visual-lexeme-spans
% Source: figures/lexeme-spans.tex:3 (tikzpicture env start line)
% Original label: \label{fig:visual-lexeme-spans}
% Original caption (LaTeX source):
%   Half-open source spans distinguish three tokens in five characters; the final cursor points just past the input.
% Original caption (plain text):
%   Half-open source spans distinguish three tokens in five characters; the final cursor points just past the input.
% Required TikZ setup: \usepackage{tikz} + \usetikzlibrary{arrows.meta, positioning, shapes.geometric, fit} (loaded by support/compiler-pseudocode.sty).
% Shared styles used in this picture: `cfg edge'.
%   Defined in support/compiler-pseudocode.sty: `basic block', `cfg edge', `back edge', `block label'.
% Usage: % Slug: visual-lexeme-spans
% Source: figures/lexeme-spans.tex:3 (tikzpicture env start line)
% Original label: \label{fig:visual-lexeme-spans}
% Original caption (LaTeX source):
%   Half-open source spans distinguish three tokens in five characters; the final cursor points just past the input.
% Original caption (plain text):
%   Half-open source spans distinguish three tokens in five characters; the final cursor points just past the input.
% Required TikZ setup: \usepackage{tikz} + \usetikzlibrary{arrows.meta, positioning, shapes.geometric, fit} (loaded by support/compiler-pseudocode.sty).
% Shared styles used in this picture: `cfg edge'.
%   Defined in support/compiler-pseudocode.sty: `basic block', `cfg edge', `back edge', `block label'.
% Usage: \input{graphs-code/visual-lexeme-spans.tikz.tex} inside a figure environment
%   (wrap with \begin{figure} ... \caption{...} \label{fig:visual-lexeme-spans} ... \end{figure}).
% Note: tikzpicture body below is verbatim from the source; do not reformat.
\begin{tikzpicture}[x=14mm,y=10mm]
    \foreach \x/\c in {0/x,1/<,2/=,3/1,4/0} {
      \draw (\x,0) rectangle +(1,1);
      \node[font=\ttfamily] at (\x+0.5,0.5) {\c};
      \node[font=\small] at (\x,1.3) {\x};
    }
    \node[font=\small] at (5,1.3) {5};
    \node[font=\sffamily\small,anchor=east] at (-0.2,1.3) {Offset};
    \draw[thick] (0,-0.2) -- (1,-0.2);
    \draw[thick] (1,-0.5) -- (3,-0.5);
    \draw[thick] (3,-0.2) -- (5,-0.2);
    \node[font=\sffamily\small,align=center] at (0.5,-1.1) {identifier\\$[0,1)$};
    \node[font=\sffamily\small,align=center] at (2,-1.1) {operator\\$[1,3)$};
    \node[font=\sffamily\small,align=center] at (4,-1.1) {number\\$[3,5)$};
    \draw[cfg edge] (5,-1.8) -- (5,-0.1);
    \node[font=\sffamily\small,anchor=east] at (4.9,-1.8) {cursor after the number};
  \end{tikzpicture}
 inside a figure environment
%   (wrap with \begin{figure} ... \caption{...} \label{fig:visual-lexeme-spans} ... \end{figure}).
% Note: tikzpicture body below is verbatim from the source; do not reformat.
\begin{tikzpicture}[x=14mm,y=10mm]
    \foreach \x/\c in {0/x,1/<,2/=,3/1,4/0} {
      \draw (\x,0) rectangle +(1,1);
      \node[font=\ttfamily] at (\x+0.5,0.5) {\c};
      \node[font=\small] at (\x,1.3) {\x};
    }
    \node[font=\small] at (5,1.3) {5};
    \node[font=\sffamily\small,anchor=east] at (-0.2,1.3) {Offset};
    \draw[thick] (0,-0.2) -- (1,-0.2);
    \draw[thick] (1,-0.5) -- (3,-0.5);
    \draw[thick] (3,-0.2) -- (5,-0.2);
    \node[font=\sffamily\small,align=center] at (0.5,-1.1) {identifier\\$[0,1)$};
    \node[font=\sffamily\small,align=center] at (2,-1.1) {operator\\$[1,3)$};
    \node[font=\sffamily\small,align=center] at (4,-1.1) {number\\$[3,5)$};
    \draw[cfg edge] (5,-1.8) -- (5,-0.1);
    \node[font=\sffamily\small,anchor=east] at (4.9,-1.8) {cursor after the number};
  \end{tikzpicture}
 inside a figure environment
%   (wrap with \begin{figure} ... \caption{...} \label{fig:visual-lexeme-spans} ... \end{figure}).
% Note: tikzpicture body below is verbatim from the source; do not reformat.
\begin{tikzpicture}[x=14mm,y=10mm]
    \foreach \x/\c in {0/x,1/<,2/=,3/1,4/0} {
      \draw (\x,0) rectangle +(1,1);
      \node[font=\ttfamily] at (\x+0.5,0.5) {\c};
      \node[font=\small] at (\x,1.3) {\x};
    }
    \node[font=\small] at (5,1.3) {5};
    \node[font=\sffamily\small,anchor=east] at (-0.2,1.3) {Offset};
    \draw[thick] (0,-0.2) -- (1,-0.2);
    \draw[thick] (1,-0.5) -- (3,-0.5);
    \draw[thick] (3,-0.2) -- (5,-0.2);
    \node[font=\sffamily\small,align=center] at (0.5,-1.1) {identifier\\$[0,1)$};
    \node[font=\sffamily\small,align=center] at (2,-1.1) {operator\\$[1,3)$};
    \node[font=\sffamily\small,align=center] at (4,-1.1) {number\\$[3,5)$};
    \draw[cfg edge] (5,-1.8) -- (5,-0.1);
    \node[font=\sffamily\small,anchor=east] at (4.9,-1.8) {cursor after the number};
  \end{tikzpicture}
 inside a figure environment
%   (wrap with \begin{figure} ... \caption{...} \label{fig:visual-lexeme-spans} ... \end{figure}).
% Note: tikzpicture body below is verbatim from the source; do not reformat.
\begin{tikzpicture}[x=14mm,y=10mm]
    \foreach \x/\c in {0/x,1/<,2/=,3/1,4/0} {
      \draw (\x,0) rectangle +(1,1);
      \node[font=\ttfamily] at (\x+0.5,0.5) {\c};
      \node[font=\small] at (\x,1.3) {\x};
    }
    \node[font=\small] at (5,1.3) {5};
    \node[font=\sffamily\small,anchor=east] at (-0.2,1.3) {Offset};
    \draw[thick] (0,-0.2) -- (1,-0.2);
    \draw[thick] (1,-0.5) -- (3,-0.5);
    \draw[thick] (3,-0.2) -- (5,-0.2);
    \node[font=\sffamily\small,align=center] at (0.5,-1.1) {identifier\\$[0,1)$};
    \node[font=\sffamily\small,align=center] at (2,-1.1) {operator\\$[1,3)$};
    \node[font=\sffamily\small,align=center] at (4,-1.1) {number\\$[3,5)$};
    \draw[cfg edge] (5,-1.8) -- (5,-0.1);
    \node[font=\sffamily\small,anchor=east] at (4.9,-1.8) {cursor after the number};
  \end{tikzpicture}
